% ACTIVE-STUDY-SOURCE: Study 5 abstract.
% Edit here for both standalone and dissertation builds.

This study examines whether respiratory phase modulates peripersonal space (PPS) boundaries, the multisensory zone surrounding the body within reaching distance. We test the prediction that inhalation extends and exhalation contracts the PPS sigmoid, a prediction consistent with embodied-cognition accounts of body-referenced multisensory integration. Study~1 uses a controlled 4--4--4--4 box-breathing paradigm with audio--tactile stimulation delivered during breath holds to estimate PPS boundaries at distinct respiratory phases. Study~2 repeats the procedure with a head-mounted display that couples a visual element of the environment to the participant's ongoing breath signal, an arrangement we refer to as exteroceptive--interoceptive sensory substitution. We predict that respiratory phase will modulate the PPS sigmoid, with inhalation associated with audio--tactile facilitation at greater source distances than exhalation, and that breath-coupled visual feedback will modulate this respiratory effect relative to a matched control without spatial coupling.
